\documentclass[11pt,addpoints,answers]{exam}

%-----------------------------------------------------------------------------
% PACKAGES AND OTHER DOCUMENT CONFIGURATIONS
%-----------------------------------------------------------------------------

\usepackage[margin=1in]{geometry}
\usepackage{amsmath, amsfonts}
\usepackage{enumerate}
\usepackage{graphicx}
\usepackage{titling}
\usepackage{url}
\usepackage{xfrac}
\usepackage{natbib}
\usepackage{amssymb}
\usepackage{amsthm}
\usepackage{paralist}
\usepackage{epstopdf}
\usepackage{tabularx}
\usepackage{longtable}
\usepackage{multirow}
\usepackage{multicol}
\usepackage[colorlinks=true,urlcolor=blue]{hyperref}
\usepackage{algorithm}
\usepackage{algorithmicx}
\usepackage[noend]{algpseudocode}
\usepackage{float}
\usepackage{enumerate}
\usepackage{array}
\usepackage{environ}
\usepackage{times}
\usepackage{textcomp}
\usepackage{caption}
\usepackage{parskip} % For NIPS style paragraphs.
\usepackage[compact]{titlesec} % Less whitespace around titles
\usepackage[inline]{enumitem} % For inline enumerate* and itemize*
\usepackage{datetime}
\usepackage{comment}
% \usepackage{minted}
\usepackage{lastpage}
\usepackage{color}
\usepackage{xcolor}
\usepackage[final]{listings}
\usepackage{tikz}
\usetikzlibrary{shapes,decorations}
\usepackage{framed}
\usepackage{booktabs}
\usepackage{cprotect}
\usepackage{verbatimbox}
\usepackage{multicol}
\usepackage{hyperref}
\usepackage{subcaption}
\usepackage{mathtools} % For drcases
\usepackage{cancel}
\usepackage[many]{tcolorbox}
\usepackage{soul}
\usepackage[bottom]{footmisc}
\usepackage{bm}
\usepackage{wasysym}

\usepackage{tikz}
\usetikzlibrary{shapes.geometric}
\usetikzlibrary{positioning, arrows, automata, calc}

\newtcolorbox[]{your_solution}[1][]{
    % breakable,
    enhanced,
    nobeforeafter,
    colback=white,
    title=Your Answer,
    sidebyside align=top,
    box align=top,
    #1
}

%%%%%%%%%%%%%%%%%%%%%%%%%%%%%%%%%%%%%%%%%%%
% Formatting for \CorrectChoice of "exam" %
%%%%%%%%%%%%%%%%%%%%%%%%%%%%%%%%%%%%%%%%%%%

\CorrectChoiceEmphasis{}
\checkedchar{\blackcircle}

%%%%%%%%%%%%%%%%%%%%%%%%%%%%%%%%%%%%%%%%%%%
% Rotated Column Headers                  %
%%%%%%%%%%%%%%%%%%%%%%%%%%%%%%%%%%%%%%%%%%%
\usepackage{adjustbox}
\usepackage{array}

%https://tex.stackexchange.com/questions/32683/rotated-column-titles-in-tabular

\newcolumntype{R}[2]{%
    >{\adjustbox{angle=#1,lap=\width-(#2)}\bgroup}%
    l%
    <{\egroup}%
}
\newcommand*\rot{\multicolumn{1}{R{45}{1em}}}% no optional argument here, please!


%%%%%%%%%%%%%%%%%%%%%%%%%%%%%%%%%%%%%%%%%%%
% Better numbering                        %
%%%%%%%%%%%%%%%%%%%%%%%%%%%%%%%%%%%%%%%%%%%

% \numberwithin{equation}{section} % Number equations within sections (i.e. 1.1, 1.2, 2.1, 2.2 instead of 1, 2, 3, 4)
% \numberwithin{figure}{section} % Number figures within sections (i.e. 1.1, 1.2, 2.1, 2.2 instead of 1, 2, 3, 4)
% \numberwithin{table}{section} % Number tables within sections (i.e. 1.1, 1.2, 2.1, 2.2 instead of 1, 2, 3, 4)

%%%%%%%%%%%%%%%%%%%%%%%%%%%%%%%%%%%%%%%%%%
% Custom commands                        %
%%%%%%%%%%%%%%%%%%%%%%%%%%%%%%%%%%%%%%%%%%

\newcommand{\vc}[1]{\boldsymbol{#1}}
\newcommand{\adj}[1]{\frac{d J}{d #1}}
\newcommand{\chain}[2]{\adj{#2} = \adj{#1}\frac{d #1}{d #2}}

\newcommand{\R}{\mathbb{R}}
\newcommand{\blackcircle}{\tikz\draw[black,fill=black] (0,0) circle (1ex);}
\renewcommand{\circle}{\tikz\draw[black] (0,0) circle (1ex);}

\newcommand{\emptysquare}{{\LARGE $\square$}\ \ }
\newcommand{\filledsquare}{{\LARGE $\boxtimes$}\ \ }
\newcommand{\emptycircle}{{\LARGE $\fullmoon$}\ \ }
\newcommand{\filledcircle}{{\LARGE $\newmoon$}\ \ }

\newcommand{\ntset}{test}

% mathcal
\newcommand{\Ac}{\mathcal{A}}
\newcommand{\Bc}{\mathcal{B}}
\newcommand{\Cc}{\mathcal{C}}
\newcommand{\Dc}{\mathcal{D}}
\newcommand{\Ec}{\mathcal{E}}
\newcommand{\Fc}{\mathcal{F}}
\newcommand{\Gc}{\mathcal{G}}
\newcommand{\Hc}{\mathcal{H}}
\newcommand{\Ic}{\mathcal{I}}
\newcommand{\Jc}{\mathcal{J}}
\newcommand{\Kc}{\mathcal{K}}
\newcommand{\Lc}{\mathcal{L}}
\newcommand{\Mc}{\mathcal{M}}
\newcommand{\Nc}{\mathcal{N}}
\newcommand{\Oc}{\mathcal{O}}
\newcommand{\Pc}{\mathcal{P}}
\newcommand{\Qc}{\mathcal{Q}}
\newcommand{\Rc}{\mathcal{R}}
\newcommand{\Sc}{\mathcal{S}}
\newcommand{\Tc}{\mathcal{T}}
\newcommand{\Uc}{\mathcal{U}}
\newcommand{\Vc}{\mathcal{V}}
\newcommand{\Wc}{\mathcal{W}}
\newcommand{\Xc}{\mathcal{X}}
\newcommand{\Yc}{\mathcal{Y}}
\newcommand{\Zc}{\mathcal{Z}}

% mathbb
\newcommand{\Ab}{\mathbb{A}}
\newcommand{\Bb}{\mathbb{B}}
\newcommand{\Cb}{\mathbb{C}}
\newcommand{\Db}{\mathbb{D}}
\newcommand{\Eb}{\mathbb{E}}
\newcommand{\Fb}{\mathbb{F}}
\newcommand{\Gb}{\mathbb{G}}
\newcommand{\Hb}{\mathbb{H}}
\newcommand{\Ib}{\mathbb{I}}
\newcommand{\Jb}{\mathbb{J}}
\newcommand{\Kb}{\mathbb{K}}
\newcommand{\Lb}{\mathbb{L}}
\newcommand{\Mb}{\mathbb{M}}
\newcommand{\Nb}{\mathbb{N}}
\newcommand{\Ob}{\mathbb{O}}
\newcommand{\Pb}{\mathbb{P}}
\newcommand{\Qb}{\mathbb{Q}}
\newcommand{\Rb}{\mathbb{R}}
\newcommand{\Sb}{\mathbb{S}}
\newcommand{\Tb}{\mathbb{T}}
\newcommand{\Ub}{\mathbb{U}}
\newcommand{\Vb}{\mathbb{V}}
\newcommand{\Wb}{\mathbb{W}}
\newcommand{\Xb}{\mathbb{X}}
\newcommand{\Yb}{\mathbb{Y}}
\newcommand{\Zb}{\mathbb{Z}}

% mathbf lowercase
\newcommand{\av}{\mathbf{a}}
\newcommand{\bv}{\mathbf{b}}
\newcommand{\cv}{\mathbf{c}}
\newcommand{\dv}{\mathbf{d}}
\newcommand{\ev}{\mathbf{e}}
\newcommand{\fv}{\mathbf{f}}
\newcommand{\gv}{\mathbf{g}}
\newcommand{\hv}{\mathbf{h}}
\newcommand{\iv}{\mathbf{i}}
\newcommand{\jv}{\mathbf{j}}
\newcommand{\kv}{\mathbf{k}}
\newcommand{\lv}{\mathbf{l}}
\newcommand{\mv}{\mathbf{m}}
\newcommand{\nv}{\mathbf{n}}
\newcommand{\ov}{\mathbf{o}}
\newcommand{\pv}{\mathbf{p}}
\newcommand{\qv}{\mathbf{q}}
\newcommand{\rv}{\mathbf{r}}
\newcommand{\sv}{\mathbf{s}}
\newcommand{\tv}{\mathbf{t}}
\newcommand{\uv}{\mathbf{u}}
\newcommand{\vv}{\mathbf{v}}
\newcommand{\wv}{\mathbf{w}}
\newcommand{\xv}{\mathbf{x}}
\newcommand{\yv}{\mathbf{y}}
\newcommand{\zv}{\mathbf{z}}

% mathbf uppercase
\newcommand{\Av}{\mathbf{A}}
\newcommand{\Bv}{\mathbf{B}}
\newcommand{\Cv}{\mathbf{C}}
\newcommand{\Dv}{\mathbf{D}}
\newcommand{\Ev}{\mathbf{E}}
\newcommand{\Fv}{\mathbf{F}}
\newcommand{\Gv}{\mathbf{G}}
\newcommand{\Hv}{\mathbf{H}}
\newcommand{\Iv}{\mathbf{I}}
\newcommand{\Jv}{\mathbf{J}}
\newcommand{\Kv}{\mathbf{K}}
\newcommand{\Lv}{\mathbf{L}}
\newcommand{\Mv}{\mathbf{M}}
\newcommand{\Nv}{\mathbf{N}}
\newcommand{\Ov}{\mathbf{O}}
\newcommand{\Pv}{\mathbf{P}}
\newcommand{\Qv}{\mathbf{Q}}
\newcommand{\Rv}{\mathbf{R}}
\newcommand{\Sv}{\mathbf{S}}
\newcommand{\Tv}{\mathbf{T}}
\newcommand{\Uv}{\mathbf{U}}
\newcommand{\Vv}{\mathbf{V}}
\newcommand{\Wv}{\mathbf{W}}
\newcommand{\Xv}{\mathbf{X}}
\newcommand{\Yv}{\mathbf{Y}}
\newcommand{\Zv}{\mathbf{Z}}

% bold greek lowercase
\newcommand{\alphav     }{\boldsymbol \alpha     }
\newcommand{\betav      }{\boldsymbol \beta      }
\newcommand{\gammav     }{\boldsymbol \gamma     }
\newcommand{\deltav     }{\boldsymbol \delta     }
\newcommand{\epsilonv   }{\boldsymbol \epsilon   }
\newcommand{\varepsilonv}{\boldsymbol \varepsilon}
\newcommand{\zetav      }{\boldsymbol \zeta      }
\newcommand{\etav       }{\boldsymbol \eta       }
\newcommand{\thetav     }{\boldsymbol \theta     }
\newcommand{\varthetav  }{\boldsymbol \vartheta  }
\newcommand{\iotav      }{\boldsymbol \iota      }
\newcommand{\kappav     }{\boldsymbol \kappa     }
\newcommand{\varkappav  }{\boldsymbol \varkappa  }
\newcommand{\lambdav    }{\boldsymbol \lambda    }
\newcommand{\muv        }{\boldsymbol \mu        }
\newcommand{\nuv        }{\boldsymbol \nu        }
\newcommand{\xiv        }{\boldsymbol \xi        }
\newcommand{\omicronv   }{\boldsymbol \omicron   }
\newcommand{\piv        }{\boldsymbol \pi        }
\newcommand{\varpiv     }{\boldsymbol \varpi     }
\newcommand{\rhov       }{\boldsymbol \rho       }
\newcommand{\varrhov    }{\boldsymbol \varrho    }
\newcommand{\sigmav     }{\boldsymbol \sigma     }
\newcommand{\varsigmav  }{\boldsymbol \varsigma  }
\newcommand{\tauv       }{\boldsymbol \tau       }
\newcommand{\upsilonv   }{\boldsymbol \upsilon   }
\newcommand{\phiv       }{\boldsymbol \phi       }
\newcommand{\varphiv    }{\boldsymbol \varphi    }
\newcommand{\chiv       }{\boldsymbol \chi       }
\newcommand{\psiv       }{\boldsymbol \psi       }
\newcommand{\omegav     }{\boldsymbol \omega     }

% bold greek uppercase
\newcommand{\Gammav     }{\boldsymbol \Gamma     }
\newcommand{\Deltav     }{\boldsymbol \Delta     }
\newcommand{\Thetav     }{\boldsymbol \Theta     }
\newcommand{\Lambdav    }{\boldsymbol \Lambda    }
\newcommand{\Xiv        }{\boldsymbol \Xi        }
\newcommand{\Piv        }{\boldsymbol \Pi        }
\newcommand{\Sigmav     }{\boldsymbol \Sigma     }
\newcommand{\Upsilonv   }{\boldsymbol \Upsilon   }
\newcommand{\Phiv       }{\boldsymbol \Phi       }
\newcommand{\Psiv       }{\boldsymbol \Psi       }
\newcommand{\Omegav     }{\boldsymbol \Omega     }

%%%%%%%%%%%%%%%%%%%%%%%%%%%%%%%%%%%%%%%%%%%
% Code highlighting with listings         %
%%%%%%%%%%%%%%%%%%%%%%%%%%%%%%%%%%%%%%%%%%%

\definecolor{bluekeywords}{rgb}{0.13,0.13,1}
\definecolor{greencomments}{rgb}{0,0.5,0}
\definecolor{redstrings}{rgb}{0.9,0,0}
\definecolor{light-gray}{gray}{0.95}

\newcommand{\MYhref}[3][blue]{\href{#2}{\color{#1}{#3}}}%

\definecolor{dkgreen}{rgb}{0,0.6,0}
\definecolor{gray}{rgb}{0.5,0.5,0.5}
\definecolor{mauve}{rgb}{0.58,0,0.82}

\lstdefinelanguage{Shell}{
  keywords={tar, cd, make},
  %keywordstyle=\color{bluekeywords}\bfseries,
  alsoletter={+},
  ndkeywords={python, py, javac, java, gcc, c, g++, cpp, .txt, octave, m, .tar},
  %ndkeywordstyle=\color{bluekeywords}\bfseries,
  identifierstyle=\color{black},
  sensitive=false,
  comment=[l]{//},
  morecomment=[s]{/*}{*/},
  commentstyle=\color{purple}\ttfamily,
  %stringstyle=\color{red}\ttfamily,
  morestring=[b]',
  morestring=[b]",
  backgroundcolor = \color{light-gray}
}

\lstset{columns=fixed, basicstyle=\ttfamily,
    backgroundcolor=\color{light-gray},xleftmargin=0.5cm,frame=tlbr,framesep=4pt,framerule=0pt}


%%%%%%%%%%%%%%%%%%%%%%%%%%%%%%%%%%%%%%%%%%%
% Custom box for highlights               %
%%%%%%%%%%%%%%%%%%%%%%%%%%%%%%%%%%%%%%%%%%%

% Define box and box title style
\tikzstyle{mybox} = [fill=blue!10, very thick,
    rectangle, rounded corners, inner sep=1em, inner ysep=1em]

% \newcommand{\notebox}[1]{
% \begin{tikzpicture}
% \node [mybox] (box){%
%     \begin{minipage}{\textwidth}
%     #1
%     \end{minipage}
% };
% \end{tikzpicture}%
% }

\NewEnviron{notebox}{
\begin{tikzpicture}
\node [mybox] (box){
    \begin{minipage}{\textwidth}
        \BODY
    \end{minipage}
};
\end{tikzpicture}
}

%%%%%%%%%%%%%%%%%%%%%%%%%%%%%%%%%%%%%%%%%%%
% Commands showing / hiding solutions     %
%%%%%%%%%%%%%%%%%%%%%%%%%%%%%%%%%%%%%%%%%%%

%% To HIDE SOLUTIONS (to post at the website for students), set this value to 0: 
\def\issoln{0}
% Some commands to allow solutions to be embedded in the assignment file.
\ifcsname issoln\endcsname \else \def\issoln{1} \fi
% Default to an empty solutions environ.
\NewEnviron{soln}{}{}
\if\issoln 1
% Otherwise, include solutions as below.
\RenewEnviron{soln}{
    \leavevmode\color{red}\ignorespaces
    % \textbf{Solution} \BODY
    \BODY
}{}
\fi

%% To HIDE TAGS set this value to 0:
\def\showtags{0}
%%%%%%%%%%%%%%%%
\ifcsname showtags\endcsname \else \def\showtags{1} \fi
% Default to an empty tags environ.
\NewEnviron{tags}{}{}
\if\showtags 1
% Otherwise, include solutions as below.
\RenewEnviron{tags}{
    \fbox{
    \leavevmode\color{blue}\ignorespaces
    \textbf{TAGS:} \texttt{\url{\BODY}}
    }
    \vspace{-.5em}
}{}
\fi

% Abhi messing around with examdoc
\qformat{\textbf{{\Large \thequestion \; \; \thequestiontitle \ (\totalpoints \ points)}} \hfill}
\renewcommand{\thequestion}{\arabic{question}}
\renewcommand{\questionlabel}{\thequestion.}

\renewcommand{\thepartno}{\arabic{partno}}
\renewcommand{\partlabel}{\thepartno.}
\renewcommand{\partshook}{\setlength{\leftmargin}{0pt}}

\renewcommand{\thesubpart}{\alph{subpart}}
\renewcommand{\subpartlabel}{(\thesubpart)}

\renewcommand{\thesubsubpart}{\roman{subsubpart}}
\renewcommand{\subsubpartlabel}{\thesubsubpart.}

% copied from stack overflow, as all good things are
\newcommand\invisiblesection[1]{%
  \refstepcounter{section}%
  \addcontentsline{toc}{section}{\protect\numberline{\thesection}#1}%
  \sectionmark{#1}}

% quite possibly the worst workaround i have made for this class
\newcommand{\sectionquestion}[1]{
\titledquestion{#1}
\invisiblesection{#1}
~\vspace{-1em}
}

%%%%%%%%%%%%%%%%%%%%%%%%%%%%%%%%%%%%%%%%%%%
% Commands for customizing the assignment %
%%%%%%%%%%%%%%%%%%%%%%%%%%%%%%%%%%%%%%%%%%%

\newcommand{\courseName}{10-301/10-601 Introduction to Machine Learning (Summer 2023)}
\newcommand{\hwName}{Programming Assignment 5: Unsupervised Learning}
\newcommand{\dueDate}{Thursday, July 20 at 11:59 PM}
\newcommand{\homeworktype}{\string programming}

\title{\textsc{\hwName}
%\thanks{Compiled on \today{} at \currenttime{}}
} % Title


\author{\courseName\\
\url{https://www.cs.cmu.edu/~hchai2/courses/10601/} \\
OUT: Thursday, July 13th \\
DUE: \dueDate{} \\ 
TAs: Alex, Andrew, Sofia, Tara, Markov, Neural the Narwhal
}

\date{}

%%%%%%%%%%%%%%%%%%%%%%%%%%%%%%%%%%%%%%%%%%%%%%%%%
% Useful commands for typesetting the questions %
%%%%%%%%%%%%%%%%%%%%%%%%%%%%%%%%%%%%%%%%%%%%%%%%%

\newcommand \expect {\mathbb{E}}
\newcommand \mle [1]{{\hat #1}^{\rm MLE}}
\newcommand \map [1]{{\hat #1}^{\rm MAP}}
\newcommand \argmax {\operatorname*{argmax}}
\newcommand \argmin {\operatorname*{argmin}}
\newcommand \code [1]{{\tt #1}}
\newcommand \datacount [1]{\#\{#1\}}
\newcommand \ind [1]{\mathbb{I}\{#1\}}

%%%%%%%%%%%%%%%%%%%%%%%%%%
% Document configuration %
%%%%%%%%%%%%%%%%%%%%%%%%%%

% Don't display a date in the title and remove the white space
\predate{}
\postdate{}
\date{}

% Don't display an author and remove the white space
%\preauthor{}
%\postauthor{}

%%%%%%%%%%%%%%%%%%
% Begin Document %
%%%%%%%%%%%%%%%%%% 

\begin{document}

\maketitle

\begin{notebox}
\paragraph{Summary} In this assignment, you will develop cover the K-Means algorithm along with its variant, K-Means++. In the programming component, you will implement an end-to-end system for both unsupervised learning algorithms. Using the MNIST dataset, you will compare the performance of the these algorithms for different values of K.
\end{notebox}
\section*{START HERE: Instructions}
\begin{itemize}
\item \textbf{Collaboration policy:} Collaboration on solving the homework is allowed, after you have thought about the problems on your own. It is also OK to get clarification (but not solutions) from books or online resources, again after you have thought about the problems on your own. There are two requirements: first, cite your collaborators fully and completely (e.g., ``Jane explained to me what is asked in Question 2.1''). Second, write your solution {\em independently}: close the book and all of your notes, and send collaborators out of the room, so that the solution comes from you only.  See the Academic Integrity Section on the course site for more information: \url{https://www.cs.cmu.edu/~hchai2/courses/10601/}

\item\textbf{Late Submission Policy:} See the late submission policy here: \url{https://www.cs.cmu.edu/~hchai2/courses/10601/}

\item\textbf{Submitting your work:} 

\begin{itemize}

% Since we are not using Canvas this semester.
% \item \textbf{Canvas:} We will use an online system called Canvas for short answer and multiple choice questions. You can log in with your Andrew ID and password. (As a reminder, never enter your Andrew password into any website unless you have first checked that the URL starts with "https://" and the domain name ends in ".cmu.edu" -- but in this case it's OK since both conditions are met).  You may only \textbf{submit once} on canvas, so be sure of your answers before you submit.  However, canvas allows you to work on your answers and then close out of the page and it will save your progress.  You will not be granted additional submissions, so please be confident of your solutions when you are submitting your assignment.

\item \textbf{Programming:} You will submit your code for programming questions on the homework to Gradescope (\url{https://gradescope.com}). After uploading your code, our grading scripts will autograde your assignment by running your program on a virtual machine (VM). \textbf{You are only permitted to use \href{https://docs.python.org/3/library/}{the Python Standard Library modules} and \texttt{numpy}}. When you are developing, check that the version number of the programming language environment (e.g. Python 3.9.12) and versions of permitted libraries (\texttt{numpy} 1.23.0) match those used on Gradescope. You have a \textbf{total of 10 Gradescope programming submissions.} Use them wisely. In order to not waste code submissions, we recommend debugging your implementation on your local machine (or the linux servers) and making sure your code is running correctly first before any Gradescope coding submission.

\item \textbf{Written:} For written problems such as short answer, multiple choice, derivations, proofs, or plots, please use the provided template. You must typeset your submission using \LaTeX{}. If your submission is misaligned with the template, there will be a \textbf{\textcolor{red}{2\% penalty}} (e.g., if the homework is out of 100 points, 2 points will be deducted from your final score). Each derivation/proof should be completed in the boxes provided. Do not move or resize any of the answer boxes. If you do not follow the template, your assignment may not be graded correctly by our AI assisted grader.

\end{itemize}

\end{itemize}

%Homework 9 will be on Gradescope, but will be "Canvas-style"- all problems will be multiple choice, select all that apply, or numerical answer. 

For multiple choice or select all that apply questions, shade in the box or circle in the template document corresponding to the correct answer(s) for each of the questions. For \LaTeX{} users, replace \lstinline{\choice} with \lstinline{\CorrectChoice} to obtain a shaded box/circle, and don't change anything else.
\clearpage

\section*{Instructions for Specific Problem Types}

For ``Select One" questions, please fill in the appropriate bubble completely:

\begin{quote}
\textbf{Select One:} Who taught this course?
    \begin{checkboxes}
     \CorrectChoice Henry Chai
     \choice Marie Curie
     \choice Noam Chomsky
    \end{checkboxes}
\end{quote}

If you need to change your answer, you may cross out the previous answer and bubble in the new answer:

\begin{quote}
\textbf{Select One:} Who taught this course?
    {
    \begin{checkboxes}
     \CorrectChoice Henry Chai
     \choice Marie Curie \checkboxchar{\xcancel{\blackcircle}{}}
     \choice Noam Chomsky
    \end{checkboxes}
    }
\end{quote}

For ``Select all that apply" questions, please fill in all appropriate squares completely:

\begin{quote}
\textbf{Select all that apply:} Which are scientists?
    {%
    \checkboxchar{$\Box$} \checkedchar{$\blacksquare$} % change checkbox style locally
    \begin{checkboxes}
    \CorrectChoice Stephen Hawking 
    \CorrectChoice Albert Einstein
    \CorrectChoice Isaac Newton
    \choice I don't know
    \end{checkboxes}
    }
\end{quote}

Again, if you need to change your answer, you may cross out the previous answer(s) and bubble in the new answer(s):

\begin{quote}
\textbf{Select all that apply:} Which are scientists?
    {%
    \checkboxchar{\xcancel{$\blacksquare$}} \checkedchar{$\blacksquare$} % change checkbox style locally
    \begin{checkboxes}
    \CorrectChoice Stephen Hawking 
    \CorrectChoice Albert Einstein
    \CorrectChoice Isaac Newton
    \choice I don't know
    \end{checkboxes}
    }
\end{quote}

For questions where you must fill in a blank, please make sure your final answer is fully included in the given space. You may cross out answers or parts of answers, but the final answer must still be within the given space.

\begin{quote}
\textbf{Fill in the blank:} What is the course number?

\begin{tcolorbox}[fit,height=1cm, width=4cm, blank, borderline={1pt}{-2pt},nobeforeafter]
    \begin{center}\huge10-601\end{center}
    \end{tcolorbox}\hspace{2cm}
    \begin{tcolorbox}[fit,height=1cm, width=4cm, blank, borderline={1pt}{-2pt},nobeforeafter]
    \begin{center}\huge10-\xcancel{6}301\end{center}
    \end{tcolorbox}
\end{quote}
\clearpage

{\LARGE \bf Written Questions (2 \ points)}
\begin{questions}
\sectionquestion{PCA}

\begin{parts}
\uplevel{
    For this section, refer to the PCA demo linked \href{https://colab.research.google.com/drive/1KF2pVCqQuQABJz_tlSNu1WeHkboV7YG_?usp=sharing}{here}. In this demonstration, we have performed PCA for you on a \href{https://en.wikipedia.org/wiki/Iris_flower_data_set}{simple four-feature dataset}. The questions below have also been added to the colab notebook linked for ease of access. Run the code in the notebook, then answer the questions based on the results.
    }
 \part[1] \sone Do you see any special relationships between any of the features? In particular, take a look at the \texttt{petal\_length} feature. How would you describe its association with each of the \textbf{other features}? Select the correct statement with appropriate justification.

    \begin{checkboxes}
    % YOUR ANSWER
    % Change \choice to \CorrectChoice for the appropriate selection/selections
    \choice The features are highly correlated: we observe linearly proportional relationships where increases in \verb|petal_length| often correspond to increases in another feature
    \choice The features are highly correlated: we observe that the color classes can be separated with decision boundaries along the \verb|petal_length| axis.
    \choice The features are uncorrelated: we observe random noise as if the features were generated from independent distributions
    \choice The features are uncorrelated: we observe the ``default $y=x$'' relationship between features
    \end{checkboxes}

      \part[1]  If we wanted to find $k$ principal components such that we preserve \textbf{at least} 95\% of the variance in the data, what would be the value of $k$? Hint: it is helpful here to look at the cumulative variance in the first $k$ components, which we have calculated for you.
    
    \begin{your_solution}[height=2cm, width=.15\textwidth, title=$k$]
    \end{your_solution}
\end{parts}
\clearpage

\sectionquestion{Empirical Questions}
\\
The following questions should be completed after you work through the programming portion of this assignment. \textbf{For any plotting questions, you must title your graph, label your axes and provide units (if applicable), and provide a legend (if applicable) in order to receive full credit.}
\\
Please submit computer-generated plots for all parts. We have provided you with the code to generate these plots in the handout.

\begin{parts}
\part K-means

The following questions should be completed after you work through the programming portion of this assignment. Use \verb|mnist_train.csv| for the following questions.

\begin{subparts}
\subpart[3] Run the K-Means algorithm with K=2. Include plots below of the cluster centers returned by the K-Means algorithm.

\begin{your_solution}[height=7cm, width=\textwidth]
\end{your_solution}


\subpart[3] Run the K-Means algorithm with K=5. Include plots below of the cluster centers returned by the K-Means algorithm.

\begin{your_solution}[height=7cm, width=\textwidth]
\end{your_solution}


\subpart[3] Run the K-Means algorithm with K=10. Include plots below of the cluster centers returned by the K-Means algorithm.

\begin{your_solution}[height=7cm, width=\textwidth]
\end{your_solution}


\subpart[2] Compare the performance of the three values of K and discuss which one is optimal and why.

\begin{your_solution}[height=3cm, width=\textwidth]
\end{your_solution}


\newpage
\subpart[3] Run the K-Means++ algorithm with K=5. Include plots below of the cluster centers returned by the K-Means++ algorithm.

\begin{your_solution}[height=7cm, width=\textwidth]
\end{your_solution}


\subpart[3] Run the K-Means++ algorithm with K=10. Include plots below of the cluster centers returned by the K-Means++ algorithm.

\begin{your_solution}[height=7cm, width=\textwidth]
\end{your_solution}


\subpart[2] Compare the performance of the K-Means ++ algorithm with the K-Means algorithm using the loss values of both with K=10 and explain which algorithm performs better than the other and why.

\begin{your_solution}[height=3cm, width=\textwidth]
\end{your_solution}

\end{subparts}
\end{parts}


\newpage
\end{questions}
\newpage
\section{Collaboration Questions}
After you have completed all other components of this assignment, report your answers to these questions regarding the collaboration policy. Details of the policy can be found \href{http://www.cs.cmu.edu/~hchai/courses/10601/syllabus.html}{here}.
\begin{enumerate}
    \item Did you receive any help whatsoever from anyone in solving this assignment? If so, include full details.
    \item Did you give any help whatsoever to anyone in solving this assignment? If so, include full details.
    \item Did you find or come across code that implements any part of this assignment? If so, include full details.
\end{enumerate}

\begin{your_solution}[height=6cm]
% YOUR ANSWER 

\end{your_solution}
\newpage

{\LARGE \bf Programming (41 points)}
\\
% Welcome to the programming component of this assignment! 

\section{The Task}
Our goal in this assignment is to implement a K-Means algorithm to cluster images in the MNIST dataset.

\section{The Datasets}
\label{sec:dataset}

\paragraph{Dataset} We will be using a \textbf{subset} of the Modified National Institute of Standards and Technology (MNIST) database. This data includes images of all 10 handwritten digits; our subset will include 3000 of the available images in the dataset. We will evaluate your code on this subset of 3000 samples.

\paragraph{File Format} The dataset consists of two csv files. One of these file uses a subset of the 3000 samples given to you. The digits in this subset are restricted to only 0 and 1. Each row contains 784 columns separated by commas. These represent the pixel values of the image.


\section{Model Definition}

In this assignment you will implement the K-means function to perform K-means clustering  on the MNIST dataset. You will also implement the K-Means variant, K-Means++ on the same dataset. The input to your dataset will be N 784-dimensional data points and your output will be the cluster assignments of each of the data points in your training input.

Your code should be able to run with K-Means or K-Means++ invariably on any dataset we pass in according to the algorithm flag passed into the command line when your code is run. Furthermore, you should be able to assign cluster centers for any positive nonzero value of K which is passed in. 

For K-Means, you will be asked to initialize the cluster centers to K random points. It is critical that in order for your code to run successfully, you do not change the random seed defined at the top of the file \texttt{kmeans.py}.  For K-Means++, you must randomly assign your first cluster center, and then use the definition of the K-Means++ algorithm to define the subsequent K-1 cluster centers.

\subsection{Command Line Arguments}
\label{sec:args}
The autograder runs and evaluates the output from the files generated, using the following command:

\begin{lstlisting}[language=Shell]
$ python3 kmeans.py [args...]
\end{lstlisting}

Where above \texttt{[args\dots]} is a placeholder for four command-line arguments: \texttt{<train\_input>}\newline \texttt{<K>} \texttt{<algorithm>} \texttt{<train\_output>}. These arguments are described in detail below:
\begin{enumerate}
    \item \texttt{<train\_input>}: path to the training input \texttt{.csv} file (see Section~\ref{sec:dataset})
    \item \texttt{<K>}: integer specifying the number of cluster centers which should be initialized.
    \item \texttt{<algorithm>}: integer taking value 0 or 1 that specifies whether to use {\sc K-Means} or {\sc K-Means++}--- that is, if \lstinline{algorithm}==0 use K-Means, and if \lstinline{algorithm}==1 use K-Means++ initialization.
    \item \texttt{<train\_output>}: cluster assignments for all data points in \texttt{train\_input}.
    
\end{enumerate}

As an example, the following command line would run your program with K=2 on the small data provided in the handout using K-Means.
\\
\begin{lstlisting}[language=Shell]
python3 kmeans.py mnist_small_train.csv 2 0 small_train_out.txt
\end{lstlisting}

The command line arguments are parsed for you in \texttt{kmeans.py} using the Python builtin \texttt{argparse} package.

\subsection{Sample Outputs}

In the handout file we have also provided you with sample outputs for the small dataset. These files correspond to the expected output for the small dataset with K=2 and K=5 as well as algorithm=0 and algorithm =1. Although you will not be required to output the final cluster centers after your k-means algorithm terminates, we have included output files containing the final cluster centers for the four settings described above to help you debug your code.

\subsection{Gradescope Submission}
You should only submit your \texttt{kmeans.py} file to Gradescope. Please do not use any other file name for your implementation. This will cause problems for the autograder to correctly detect and run your code.

\textbf{Note:} For this assignment, you may make up to \textbf{10} submissions to Gradescope before the deadline, but only your last submission will be graded.
\newpage


% \appendix
% \input{appendix.tex}

\end{document}

